\section*{Программа экзамена}
\addcontentsline{toc}{section}{Программа экзамена} % Добавляем в оглавление

\begin{enumerate}[leftmargin=*, wide=0pt]
    \item Ранг матрицы. Теорема о базисном миноре. Теорема о ранге матрицы.
    
    \item Системы линейных уравнений. Метод Гаусса. Теорема Кронекера-Капелли. Фундаментальная система решений и общее решение однородной системы линейных уравнений. Общее решение неоднородной системы. Теорема Фредгольма.
    
    \item Аксиоматика линейного пространства. Линейная зависимость и линейная независимость систем элементов в линейном пространстве. Базис и размерность.
    
    \item Координатное представление векторов линейного пространства и операций с ними. Теорема об изоморфизме. Матрица перехода от одного базиса к другому. Изменение координат при изменении базиса в линейном пространстве.
    
    \item Подпространства и способы их задания в линейном пространстве. Сумма и пересечение подпространств. Формула размерности суммы подпространств. Прямая сумма.
    
    \item Линейные отображения линейных пространств и линейные преобразования линейного пространства. Ядро и образ линейного отображения. Операции над линейными преобразованиями. Обратное преобразование. Линейное пространство линейных отображений (преобразований).
    
    \item Матрицы линейного отображения и линейного преобразования для конечномерных пространств. Операции над линейными преобразованиями в матричной форме. Изменение матрицы линейного отображения (преобразования) при замене базисов. Изоморфизм пространства линейных отображений и пространства матриц.
    
    \item Инвариантные подпространства линейных преобразований. Собственные векторы и собственные значения. Собственные подпространства. Линейная независимость собственных векторов, принадлежащих различным собственным значениям.
    
    \item Нахождение собственных значений и собственных векторов линейного преобразования конечномерного линейного пространства. Характеристическое уравнение, его инвариантность. Оценка размерности собственного подпространства. Условия диагонализуемости матрицы линейного преобразования. Теорема Гамильтона-Кэли.
    
    \item Линейные формы. Сопряженное (двойственное) пространство. Биортогональный базис.
    
    \item Билинейные и квадратичные формы. Их координатное представление в конечномерном линейном пространстве. Изменение матриц билинейной и квадратичной форм при изменении базиса.
    
    \item Приведение квадратичной формы к каноническому виду методом Лагранжа. Теорема (закон) инерции для квадратичных форм. Знакоопределенные квадратичные формы. Критерий Сильвестра. Приведение квадратичной формы к каноническому виду элементарными преобразованиями.
    
    \item Аксиоматика евклидова пространства. Неравенство Коши-Буняковского. Неравенство треугольника. Матрица Грама и ее свойства.
    
    \item Процесс ортогонализации в евклидовом пространстве. Переход от одного ортонормированного базиса к другому. Ортогональное дополнение подпространства, ортогональное проектирование на подпространство.
    
    \item Линейные преобразования евклидова пространства. Сопряженные преобразования, их свойства. Матрица сопряженного преобразования.
    
    \item Самосопряженные преобразования. Свойства их собственных векторов и собственных значений. Существование ортонормированного базиса из собственных векторов самосопряженного преобразования. Ортогональное проектирование на подпространство как пример самосопряженного преобразования.
    
    \item Ортогональные преобразования. Их свойства. Ортогональные матрицы.
    
    \item Полярное разложение линейных преобразований евклидова пространства. Сингулярное разложение.
    
    \item Построение ортонормированного базиса, в котором квадратичная форма имеет диагональный вид. Одновременное приведение к диагональному виду пары квадратичных форм, одна из которых является знакоопределенной.
\end{enumerate}